\section{How many shifted Jensen tests detect a polynomial perturbation?}
\label{sec:finite-jensen}
The incoming manuscript \href{https://github.com/VladimirReshetnikov/ProveIt/blob/5cc267af75a635199ac5d93a7d49cbce521daefc/docs/incoming/grh_transfers_density_gaps_and_zero_tests.zip}{\emph{Geometric progressions and nonabelian zero tests}} (files \texttt{jensen.tex} and \texttt{gaussian.tex})
constructs cubic counterexamples to a general kernel criterion while studying
the conductor-$23$ Artin factor. Its Gaussian argument uses every shift of a
cubic Jensen polynomial. For polynomial perturbations there is an exact finite
answer. This calculation determines what evidence from small perturbations
can establish. It does not assume any Riemann hypothesis.

For a real polynomial $V$ bounded below on $[0,\infty)$, consider the even kernel
$K_\epsilon(u)=\exp[-u^2-\epsilon V(u^2)]$, $\epsilon>0$.
Define the coefficients $\gamma_n$ by
\[
 \int_{\mathbb R} e^{zu}K_\epsilon(u)\,du
 =\sum_{n\ge0}\gamma_n z^{2n}/n!,\qquad
 J^{d,n}(X)=\sum_{j=0}^d \binom{d}{j}\gamma_{n+j}X^j .
\]
Call shift $n$ persistent when $J^{3,n}$ is real-rooted along some sequence
$\epsilon\downarrow0$. Each shift may use its own sequence.

\paragraph{Theorem.}
Fix $r\ge3$. For potentials of degree at most $r$, persistence at any
$r-2$ distinct nonnegative integer shifts forces $\deg V\le2$.
In this case all $J^{d,n}$ have negative real roots for every $\epsilon>0$,
$d\ge1$ and $n\ge0$. The test count is sharp. For any set $S$ of $r-3$
nonnegative integers there is a strictly convex increasing polynomial of degree
$r$ whose persistent shifts are precisely $S$. At each fixed shift, the predicted
outcome holds for all sufficiently small positive $\epsilon$.

\paragraph{Proof.}
Set $p=n+1/2$ and let $X_n$ have the gamma distribution of shape $p$
and rate one. The substitution $x=u^2$ and the gamma duplication
formula give $c_n=4^n\gamma_n/\sqrt\pi=\mathbb E e^{-\epsilon V(X_n)}$.
A cubic built from the $c_n$ differs from the original only by positive
rescalings of variable and value. Write $m(n)=\mathbb EV(X_n)$ and
$\Delta m(n)=m(n+1)-m(n)$. The Taylor estimate
$c_n=1-\epsilon m(n)+O_n(\epsilon^2)$ follows by integrating the
pointwise second-order remainder, bounded by a constant times
$\epsilon^2 V^2$ for $0<\epsilon<1$ using the lower bound on $V$.
For the translated cubic $P(W)=\sum_{j=0}^3\binom{3}{j}c_{n+j}(W-1)^j$
the coefficients of $1,W,W^2,W^3$ are, respectively,
\[
 \epsilon\Delta^3m(n)+O_n(\epsilon^2),\quad
 -3\epsilon\Delta^2m(n+1)+O_n(\epsilon^2),\quad
 O_n(\epsilon),\quad 1+O_n(\epsilon).
\]
Thus, if $\Delta^3m(n)\ne0$, its discriminant equals
$-27\epsilon^2(\Delta^3m(n))^2+O_n(\epsilon^3)$.
If $\Delta^3m(n)=0$, it equals
$108\epsilon^3(\Delta^2m(n+1))^3+O_n(\epsilon^4)$.
The transform $V\mapsto m$ takes $x^k$ to the monic polynomial
$(n+1/2)(n+3/2)\cdots(n+k-1/2)$, so is a degree-preserving
linear isomorphism. Hence $r-2$ persistent shifts force the polynomial
$\Delta^3m$, of degree at most $r-3$, to vanish identically.

For the converse take $V(x)=ax^2+bx+c$. For $a>0$, P\'olya's quartic
Fourier integral theorem with arbitrary quadratic coefficient gives only real
zeros for $\int e^{itu}K_\epsilon(u)\,du$.
See N.~G.~de Bruijn, \emph{The roots of trigonometric integrals},
Duke Math.~J.~17 (1950), 197--226,
\href{https://repository.tudelft.nl/file/File_524a4d3c-0858-4c55-b425-8c8ce1fb89da}{p.~197, (1.3)}.
The function $G(w)=\int\cosh(\sqrt w\,u)K_\epsilon(u)\,du$ then
has negative real zeros only. The bound
$|\sqrt w||u|\le (\epsilon a/2)u^4+C|w|^{2/3}$ gives order
at most $2/3$ for $G$. Its genus-zero product places it in the
Laguerre--P\'olya class. For $a=0$, boundedness below gives $b\ge0$
and $G$ is a positive multiple of an exponential with positive slope,
also in that class. Every derivative of $G$ is in the class, so Jensen's
theorem gives real roots for every $J^{d,n}$. Positivity of all the
coefficients puts these roots on the negative real axis.

Finally take $m$ with $\Delta^3m(n)=\prod_{s\in S}(n-s)$, solving by
triangularity of finite differences. Its leading coefficient and that of
the corresponding $V$ are positive. Adding $A x^2$ to $V$ adds $2A$ to
$\Delta^2m(n+1)$ and preserves $\Delta^3m$. Choose $A$ large enough for
$\Delta^2m(s+1)>0$ at each $s\in S$ and $V''>0$ throughout $[0,\infty)$.
The latter is possible since $V''$ has positive leading coefficient.
A further linear adjustment makes $V'(0)>0$. The two discriminant expansions
give all claims. The small-$\epsilon$ thresholds can depend on the shift.
\hfill$\square$

\paragraph{An exact example.}
Take $r=4$ and $V(x)=x^4-14x^3+100x^2+x$. Since
$V''(x)=12(x-7/2)^2+53$, the potential increases from zero and is strictly
convex. Its moment polynomial has $\Delta^3m(n)=24n$ and
$\Delta^2m(1)=95$. Shift zero therefore persists, and each other fixed
shift fails for sufficiently small $\epsilon$. At $\epsilon=10^{-12}$
the checker \texttt{check\_finite.py} bounds the discriminants of the
rescaled cubics using the pointwise inequalities
$1-y\le e^{-y}\le1-y+y^2/2$ for $y\ge0$. Its rational bounds are
\[
 92{,}590{,}000<\operatorname{disc}(P_0)/\epsilon^3<92{,}600{,}000,
 \qquad
 -15{,}553<\operatorname{disc}(P_1)/\epsilon^2<-15{,}551 .
\]
